\begin{proof}
\textbf{Theorem (Bernoulli’s Inequality).}  
Let \(x\in\mathbb{R}\) satisfy \(x\ge -1\) and let \(n\in\mathbb{N}_{0}\). Then  

\[
(1+x)^{n}\;\ge\;1+nx .
\]

\begin{enumerate}
\item \textbf{Degenerate cases.}
\begin{itemize}
\item \(\mathbf{n=0}\).  \((1+x)^{0}=1\) and \(1+0\cdot x=1\); equality holds.
\item \(\mathbf{x=-1}\).  For \(n\ge1\), \((1+x)^{n}=0\) and \(1+nx=1-n\le0\); hence \(0\ge1-n\). Equality occurs only when \(n=1\) (already covered by the case \(n=1\)).
\item \(\mathbf{x=0}\).  \((1+0)^{n}=1\) and \(1+n\cdot0=1\); equality holds for every \(n\).
\end{itemize}
Thus the theorem is true in all degenerate situations.

\item \textbf{Induction on the exponent.}
\begin{enumerate}
\item \textbf{Base step (\(n=1\)).}
\[
(1+x)^{1}=1+x=1+1\cdot x ,
\]
so the inequality holds with equality.

\item \textbf{Inductive hypothesis.}  
Assume that for a fixed integer \(k\ge1\)

\[
\boxed{(1+x)^{k}\;\ge\;1+kx}\qquad\text{for every }x\ge-1 . \tag{H\(_k\)}
\]

\item \textbf{Inductive step.}  
Since \(x\ge-1\) we have \(1+x\ge0\). Multiplying both sides of \((\text{H}_k)\) by this non‑negative factor preserves the inequality:

\begin{align*}
(1+x)^{k+1}
   &= (1+x)\,(1+x)^{k} \\
   &\ge (1+x)(1+kx). \tag{1}
\end{align*}

Expanding the product gives the identity  

\[
(1+x)(1+kx)=1+(k+1)x+kx^{2}. \tag{2}
\]

Because \(k\ge1\) and \(x^{2}\ge0\), the term \(kx^{2}\) is non‑negative; discarding it yields  

\[
1+(k+1)x+kx^{2}\;\ge\;1+(k+1)x . \tag{3}
\]

Combining (1), (2) and (3) we obtain  

\[
(1+x)^{k+1}\;\ge\;1+(k+1)x .
\]

Thus the statement holds for \(n=k+1\).
\end{enumerate}
By the principle of mathematical induction, \((\text{H}_n)\) is true for every integer \(n\ge0\); i.e.

\[
(1+x)^{n}\;\ge\;1+nx \qquad\forall\,x\ge-1,\; n\in\mathbb{N}_{0}. \tag{*}
\]

\item \textbf{Equality conditions.}  
Equality can arise only in the steps where we did not introduce a strict inequality:

\begin{itemize}
\item In the base cases \(n=0\) and \(n=1\) we have exact equality for all admissible \(x\).
\item In the inductive step the only possible loss of equality is the removal of the non‑negative term \(kx^{2}\) in \((3)\). Equality there forces \(kx^{2}=0\); since \(k\ge1\) this yields \(x^{2}=0\), i.e. \(x=0\).
\end{itemize}

Hence equality holds \emph{iff} one of the following occurs:
\begin{enumerate}
\item \(n=0\) (any \(x\ge-1\));
\item \(n=1\) (any \(x\ge-1\));
\item \(x=0\) (any \(n\ge0\)).
\end{enumerate}
No other combination gives equality, because for \(n\ge2\) and \(x\neq0\) the term \(kx^{2}>0\) appears at the first inductive step and yields a strict inequality that persists for larger \(n\).

\item \textbf{Conclusion.}  
For every real number \(x\ge-1\) and every integer \(n\ge0\),

\[
\boxed{(1+x)^{n}\;\ge\;1+nx},
\]

with equality precisely in the cases listed above. \(\square\)
\end{enumerate}
\end{proof}